\section{The terminal conditions and the common pair}\label{sec:common-tuning}

We give the continuation and tuning arguments of
\cite[\S\S16.2--16.4]{QED3Master}, based on
\cite[Chapter~28]{ChamberRG}.

\subsection{Continuation within the mass range}
\label{subsec:rgfl-continuation}

Write the mass and activity bounds as
\[
 w_j^{\mathrm m}=e_j^{3/4-8\varepsilon},\qquad
 w_j^{\mathrm E}=e_j^{1/4-7\varepsilon}.
\]
For a real pair $(m_{\mathrm{in}},\varepsilon_{\mathrm{in}}^{\mathrm v})$,
the initial vacuum is $\varepsilon_{\mathrm{in}}^{\mathrm v}$,
the initial member of the mass sequence is $m_{\mathrm{in}}$,
and $E_0=\mathscr E_0=\mathfrak P_0=0$.
The variable-index notation $m_j$ includes this initial member.
The recursion is
\begin{align}
 m_{j+1}&=L(m_j+m(E_j)),&
 \varepsilon_{j+1}^{\mathrm v}
       &=L^3(\varepsilon_j^{\mathrm v}
                 +\varepsilon_{\mathrm v}(E_j)+\delta\varepsilon_j^0),
                                      \label{eq:rgfl-coupling-recursion}\\
 E_{j+1}&=\mathcal L(RE_j)^{\mathrm{nat}}
                          +\mathcal LE_j^{\mathrm{det}}+E_j^*(m_j,E_j),
                                      \label{eq:rgfl-activity-recursion}\\
 \mathscr E_{j+1}
       &=\mathcal L(R\mathscr E_j)^{\mathrm{nat}}
               +\mathcal L\delta E_j^{\mathrm{det}}
                                      +E_j^*(m_j,\mathscr E_j).
                                      \label{eq:rgfl-slot-recursion}
\end{align}
The scalar family $\mathfrak P_j$ follows
\eqref{eq:rgnl-external-scalar-map}.
In the last line the correction has the external slot;
in \eqref{eq:rgfl-activity-recursion} it has none.
These are the corresponding maps proved earlier.
The trajectory exists up to a level if every preceding
$(m_j,\mathscr E_j)$ lies in $\operatorname{Ball}_j^x(4)$.
It is controlled when it exists to $k_*(R)$ and
$|m_j|\le w_j^{\mathrm m}$ throughout, including the stop.

The normalized estimates used below are
\begin{equation}\label{eq:rgfl-window-estimates}
 \begin{gathered}
 L^{-1}\frac{w_{j+1}^{\mathrm m}}{w_j^{\mathrm m}}
                  =L^{-5/8-4\varepsilon}\le\tfrac14,\qquad
 c_{\mathrm{fl}}e_j^{1/2}
          \frac{w_j^{\mathrm E}}{w_j^{\mathrm m}}
                           =c_{\mathrm{fl}}e_j^\varepsilon,\\
 c_{\mathrm{fl}}L^{-1/4+2\varepsilon}
          \frac{w_j^{\mathrm E}}{w_{j+1}^{\mathrm E}}
                     =c_{\mathrm{fl}}L^{-3/8+11\varepsilon/2}
                                                        \le\tfrac1{16},\\
 \frac{e_{j+1}^{1/4-2\varepsilon}}{w_{j+1}^{\mathrm E}}
                                      \le e_{j+1}^\varepsilon,\qquad
 \frac{e_j^{1/4-6\varepsilon}}{w_{j+1}^{\mathrm E}}
             =L^{-1/8+3\varepsilon}e_{j+1}^\varepsilon
                                      \le e_{j+1}^\varepsilon .
 \end{gathered}
\end{equation}
Their first and third bounds use the fixed inequalities
stated above; their last two use $e_{j+1}<1$ and
$\varepsilon\le1/200$.
The charge condition also gives
\begin{equation}\label{eq:rgfl-summable-charge}
 \sum_{j=0}^{k_*}(c_{\mathrm{fl}}+2)e_j^\varepsilon
             \le(c_{\mathrm{fl}}+2)e_*^\varepsilon
                          (1-L^{-\varepsilon/2})^{-1}
                                                        \le\tfrac1{16}.
\end{equation}
In particular
$c_{\mathrm{fl}}e_j^\varepsilon\le1/16$,
$2e_j^\varepsilon\le1/16$, and
$e_j\le e_j^\varepsilon\le1/48$.

\begin{proposition}\label{prop:rgfl-continuation}
If the recursively produced masses satisfy
$|m_i|\le w_i^{\mathrm m}$ for $i<j\le k_*(R)$,
then the trajectory is defined to $j$ and
\begin{equation}\label{eq:rgfl-activity-bound}
 \|E_j\|_j\le\|\mathscr E_j\|_j^x
       \le\frac{32}{15}e_j^\varepsilon w_j^{\mathrm E}
       \le4e_j^{1/4-6\varepsilon}\le w_j^{\mathrm E}.
\end{equation}
The data satisfy
$E_j=\mathscr E_j|_{\mathscr C=0}$,
$\mathfrak P_j(X,0)=0$, and belong to their real
mass, vacuum and activity spaces.
\end{proposition}

\begin{proof}
At zero all activities and scalar families vanish.
Assume the assertion through $i<j$.
The mass condition and
$4e_i^\varepsilon\le1$ put the input in the unit ball,
so Theorem~\extref{A}{NonlinearStep} applies at $i$.
The scaled connection pair lies in
$\widetilde R_i^-\times\mathfrak X_i$, so every local
determinant and correction in the recursion is defined.
The scalar external sum is finite and uses
$\mathscr C_L\in\mathfrak X_i$.
The step theorem preserves the real spaces
and commutes with restriction to $\mathscr C=0$.
Consequently the coupling maps computed from
$E_i$ or $\mathscr E_i$ coincide.
The scalar recursion vanishes at zero because each
$E_i^{\mathrm{det}}(X,0)$ vanishes.

The activity bound of the current step, divided by
$w_{i+1}^{\mathrm E}$ and combined with
\eqref{eq:rgfl-window-estimates}, gives
\[
 \frac{\|\mathscr E_{i+1}\|_{i+1}^x}{w_{i+1}^{\mathrm E}}
      \le\frac1{16}
                   \frac{\|\mathscr E_i\|_i^x}{w_i^{\mathrm E}}
                                     +2e_{i+1}^\varepsilon.
\]
Starting from zero, the bound
$(32/15)e_i^\varepsilon$ is preserved since
$e_i\le e_{i+1}$ and $(32/15)/16+2=32/15$.
Restriction does not increase the activity norm.
This proves \eqref{eq:rgfl-activity-bound} and the induction.

At the first step the formulas use $192\ell_0^3$
incoming pairs with weight $1/48$ and
$4\ell_1^3$ outgoing pairs, as proved in the
exact density identity.
Its coupling maps are
$m_1=Lm_{\mathrm{in}}$ and
$\varepsilon_1^{\mathrm v}
 =L^3(\varepsilon_{\mathrm{in}}^{\mathrm v}+\delta\varepsilon_0^0)$;
its activity is the local determinant term plus
the first correction. The normalized estimate has
zero incoming-activity term and the same two
source bounds. Thus the induction includes that step.
\end{proof}

\subsection{The terminal conditions and the fixed point}
\label{subsec:rgfl-fixed-point}

\begin{theorem}\label{thm:rgfl-terminal-zero-pair}
On any $R\in\mathcal T(e)$ with $N\ge N_{\min}(R,e)$,
there is exactly one real controlled pair whose
trajectory has $m_{k_*}=0$ and
$\varepsilon_{k_*}^{\mathrm v}=0$.
Along its trajectory,
\begin{equation}\label{eq:rgfl-tuned-bounds}
 |m_j|\le\tfrac18w_j^{\mathrm m},\qquad
 \|E_j\|_j\le\|\mathscr E_j\|_j^x
                                  \le4e_j^{1/4-6\varepsilon},\qquad
 |\varepsilon_j^{\mathrm v}|
                          \le5c_{\mathrm{fl}}e_j^{1/4-6\varepsilon}.
\end{equation}
The vacuum at every level is
\begin{equation}\label{eq:rgfl-backward-vacuum}
 \varepsilon_j^{\mathrm v}
     =-\sum_{i=j}^{k_*-1}L^{-3(i-j)}
                      \big(\varepsilon_{\mathrm v}(E_i)
                                              +\delta\varepsilon_i^0\big).
\end{equation}
\end{theorem}

\begin{proof}
Take the finite product of the real mass coordinates
and real activity Banach spaces, with norm
\[
 \|\mathbf z\|_{\mathrm{fl}}
   =\max_{0\le j\le k_*}
       \max\left(\frac{|m_j|}{w_j^{\mathrm m}},
                              \frac{\|E_j\|_j}{w_j^{\mathrm E}}\right),
                  \qquad \mathbf z=(m_j,E_j)_{0\le j\le k_*}.
\]
The closed subset $\mathcal B$ defined by these
quotients being at most $1$ and $E_0=0$
is nonempty and complete. Define a self-map by
\begin{align}
 \widehat m_j&=L^{-1}m_{j+1}-m(E_j)\quad(j<k_*),
                 &\widehat m_{k_*}&=0, \nonumber\\
 \widehat E_0&=0,\qquad
 \widehat E_{j+1}
       =\mathcal L(RE_j)^{\mathrm{nat}}
                    +\mathcal LE_j^{\mathrm{det}}+E_j^*(m_j,E_j).
                                      \label{eq:rgfl-fixed-point-map}
\end{align}
All maps here are the actual maps without an external
slot; their inputs belong to the unit ball.
Extraction, adjustment followed by reblocking,
the determinant activity and the nonlinear correction
are defined on that ball and preserve the real space.

We verify that the map preserves $\mathcal B$.
The extraction estimate implies
$|m(E_j)|\le c_{\mathrm{fl}}e_j^{1/2}\|E_j\|_j$.
Equations~\eqref{eq:rgfl-window-estimates} and
\eqref{eq:rgfl-summable-charge} give
\[
 \frac{|\widehat m_j|}{w_j^{\mathrm m}}
                    \le\tfrac14+c_{\mathrm{fl}}e_j^\varepsilon
                    \le\tfrac5{16}.
\]
For the activity, the contraction, determinant and
correction estimates give
\[
 \frac{\|\widehat E_{j+1}\|_{j+1}}{w_{j+1}^{\mathrm E}}
                         \le\tfrac1{16}+2e_{j+1}^\varepsilon
                         \le\tfrac18 .
\]
The fixed end components are zero. Hence the map
takes $\mathcal B$ into itself.

For two sequences, let $d_j^{\mathrm m},d_j^{\mathrm E}$
be the corresponding normalized differences.
Linearity of extraction gives
\[
 \widehat d_j^{\mathrm m}
      \le\tfrac14d_{j+1}^{\mathrm m}
                            +c_{\mathrm{fl}}e_j^\varepsilon d_j^{\mathrm E}
      \le\tfrac5{16}\|\mathbf z-\widetilde{\mathbf z}\|_{\mathrm{fl}} .
\]
Adjustment followed by reblocking is linear in
the activity at fixed connection.
The local determinant term is independent of the
sequence. The correction's coupling estimate
\exteq{93}{NonlinearStep} therefore gives
\[
 \widehat d_{j+1}^{\mathrm E}
      \le\tfrac1{16}d_j^{\mathrm E}
                         +e_{j+1}^\varepsilon
                                      (d_j^{\mathrm m}+d_j^{\mathrm E})
      \le\tfrac18\|\mathbf z-\widetilde{\mathbf z}\|_{\mathrm{fl}} .
\]
The whole map contracts by $5/16$.
Banach's fixed point theorem
\cite[Theorem~8]{BrightMetric2023} gives a unique
fixed point in $\mathcal B$.
Its equations are precisely the mass and activity
recursions with $m_{k_*}=0$.

Let $m_{\mathrm{in}}$ be its initial mass coordinate
and choose
\[
 \varepsilon_{\mathrm{in}}^{\mathrm v}
      =-\sum_{i=0}^{k_*-1}L^{-3i}
                     \big(\varepsilon_{\mathrm v}(E_i)
                                           +\delta\varepsilon_i^0\big).
\]
This pair is real. Its full trajectory, including
the external slot, agrees with the fixed-point mass
and activity at every level: agreement at zero is
the initial condition; at the next level
\cref{prop:rgfl-continuation} makes the
full step defined, and the fixed-point equations
identify its mass and restricted activity.
The induction reaches $k_*$ and proves control.
Solving the scalar vacuum recursion backwards gives
\[
 \varepsilon_j^{\mathrm v}
 =L^{-3(k_*-j)}\varepsilon_{k_*}^{\mathrm v}
     -\sum_{i=j}^{k_*-1}L^{-3(i-j)}
                     \big(\varepsilon_{\mathrm v}(E_i)
                                           +\delta\varepsilon_i^0\big).
\]
At $j=0$ its comparison with the chosen input forces
$\varepsilon_{k_*}^{\mathrm v}=0$, giving
\eqref{eq:rgfl-backward-vacuum}.
Conversely any controlled pair with the two terminal
zeros has its mass/activity sequence in $\mathcal B$
by the continuation proposition. It is the same fixed
point. Its initial vacuum is then the same scalar
sum. This proves uniqueness of the pair.

For the sharper mass bound put
$B_j=|m_j|/w_j^{\mathrm m}$.
The backward mass equation and the unit activity
bound imply $B_j\le B_{j+1}/4+1/16$, while
$B_{k_*}=0$.
Summing gives $B_j\le(1/16)\sum_{s\ge0}4^{-s}
=1/12\le1/8$.
The activity bound in \eqref{eq:rgfl-tuned-bounds}
is the continuation estimate.

It remains to bound the vacuum sum.
Extraction gives
$|\varepsilon_{\mathrm v}(E_i)|
\le4c_{\mathrm{fl}}e_i^{1/4-6\varepsilon}$.
For $i\ge j$ the associated summand obeys
\[
 L^{-3(i-j)}e_i^{1/4-6\varepsilon}
       =e_j^{1/4-6\varepsilon}
                         L^{-(i-j)(23/8+3\varepsilon)}
       \le e_j^{1/4-6\varepsilon}16^{-(i-j)}.
\]
Thus its sum is at most
$(64/15)c_{\mathrm{fl}}e_j^{1/4-6\varepsilon}$.
The box normalization satisfies
$0\le\delta\varepsilon_i^0\le e_i^7$ by the
step estimate.
Its remaining sum is bounded by
\[
 \sum_{i=j}^{k_*-1}L^{-3(i-j)}e_i^7
        =e_j^6\sum_{i=j}^{k_*-1}e_i
        \le e_j^6 e_*(\sqrt L-1)^{-1}
        \le\tfrac1{48}e_j^{1/4-6\varepsilon}.
\]
Here $\sqrt L\ge2$ and $e_*\le1/48$.
Since $c_{\mathrm{fl}}\ge1$,
$64/15+1/48=343/80<5$, giving the claimed bound.
The sums are empty at the stop.

At zero $E_0=0$. The backward mass equation reads
$m_{\mathrm{in}}=L^{-1}m_1$, and the forward activity is
$\mathcal LE_0^{\mathrm{det}}+E_0^*(m_{\mathrm{in}},0)$.
The determinant and correction estimates used above
include their full chamber trace and the first
weighted fluctuation integral. The initial vacuum
sum contains $\delta\varepsilon_0^0$.
No additional coupling is selected by the fixed point.
\end{proof}

\subsection{One pair on all tori of the construction}
\label{subsec:rgfl-volume-comparison}

For $N\ge N_{\min}(R_{\mathrm{top}},e)$ define
\[
 P_N(e)=\big(m_{0,N}(e),\varepsilon_{0,N}^{\mathrm v}(e)\big)
\]
to be the pair of \cref{thm:rgfl-terminal-zero-pair}
on the prescribed side $R_{\mathrm{top}}(e)$.
The dependence on the fixed $\bar m$ and fixed
construction parameters is suppressed.
For lower cutoffs set $m_{0,N}(e)=0$.

\begin{theorem}\label{thm:rgfl-common-pair}
The pair $P_N(e)$ is controlled on every
$R\in\mathcal T(e)$ with $N\ge N_{\min}(R,e)$.
Up to its stop, its mass and vacuum sequence agree
with those on the top torus; its activity and
external scalar families agree there on flat polymers.
For every such level,
\begin{equation}\label{eq:rgfl-common-bounds}
 |m_j|\le\tfrac18e_j^{3/4-8\varepsilon},\qquad
 \|E_j\|_j\le\|\mathscr E_j\|_j^x
                             \le4e_j^{1/4-6\varepsilon},\qquad
 |\varepsilon_j^{\mathrm v}|
                         \le5c_{\mathrm{fl}}e_j^{1/4-6\varepsilon}.
\end{equation}
The real mass sequence is independent of the side
datum and of the twist. At all $N\ge0$,
\begin{equation}\label{eq:rgfl-bare-mass-bound}
 |m_{0,N}(e)|\le\tfrac18q^{3/4-8\varepsilon}<\tfrac18,
                       \qquad q=eL^{-N/2}.
\end{equation}
The total lattice mass is
$\boldsymbol\mu_N=\bar m\eta+m_{0,N}(e)$.
\end{theorem}

\begin{proof}
The pair is defined because
$N_{\min}(R_{\mathrm{top}},e)\le N_{\min}(R,e)$.
The two trajectories have the same initial
couplings and zero initial activities.
Suppose they agree through $j<k_*(R)$.
The top mass is bounded by $w_j^{\mathrm m}/8$;
the same bound on the smaller side permits continuation
to $j+1$ by \cref{prop:rgfl-continuation}.
Both inputs are now in the unit balls, are real,
and have locally agreeing activity and scalar families.
The scale proposition verifies their primitive,
domain and next-side comparison hypotheses.
\cref{thm:rgcov-step-comparison} gives the
same extracted mass and vacuum, including the same
$\delta\varepsilon_j^0$, and local agreement of the
outgoing activity and scalar family.
The scalar recursions therefore give the same
$m_{j+1},\varepsilon_{j+1}^{\mathrm v}$.
Restriction to zero external slot preserves agreement
of $E_{j+1}$.
This proves the induction up to $k_*(R)$.

The common mass bound follows from the top trajectory.
The local comparison alone is not an estimate of
the entire activity on the smaller torus; that
bound follows from its own continuation estimate.
The vacuum is a common scalar and satisfies the
top bound at this same level. Thus
\eqref{eq:rgfl-common-bounds} follows.
At a smaller-volume stop the couplings need not
vanish; the zeros defining $P_N(e)$ are imposed
at $k_{\mathrm{top}}$.

The fixed-point map on the top torus is determined
by $(N,e,\bar m)$, the fixed construction parameters
and photon coordinate basis, and the top primitive
families constructed from them.
The selected top side is itself a function of
$e,\bar m$ and these parameters.
The smaller side $R$ enters no part of this definition.
Nor does a twist: the twist is a harmonic part
of a connection inserted in the density identity,
whereas the coupling maps are maps of entire
local activity families and contain no such insertion.
The full forward activity recursion retains the external slot
as a variable, whereas the fixed-point map uses its restriction
to zero. No external connection or probe is inserted when
selecting the pair.
This proves the asserted independence within
the fixed construction.

For a cutoff on which the pair is defined,
\eqref{eq:rgfl-bare-mass-bound} is its initial mass
bound on the top torus. At the other finitely many
cutoffs it holds because the assigned mass is zero.
Finally $h\ge1$ implies $e\le\bar e_1\le1/3$,
so $q\le1/3$ and its positive power is below $1$.
The sequence is real in both cases.
The sign of $\boldsymbol\mu_N$ is not fixed by this
estimate; the lattice construction admits a real mass.
\end{proof}

On the reference torus, for $N\ge N_{\mathrm{ref}}(e)$,
the outputs needed at the terminal scale are
\begin{equation}\label{eq:rgfl-reference-terminal-input}
 a_*\in\{h/L,h/L^2\},\qquad
 |m_*|\le\tfrac18e_*^{3/4-8\varepsilon},\qquad
 \|E_*\|_*\le\|\mathscr E_*\|_*^x
                                 \le4e_*^{1/4-6\varepsilon}.
\end{equation}
The primitive estimates along this flow have the
same fixed constants on every torus and at every
cutoff. The remaining terminal density and partition
estimates concern the integration of these fields.



\begin{proof}[Proof of \Cref{thm:main}]
\Cref{prop:rgfl-scales} gives the stopping levels and thresholds;
\eqref{eq:rgfl-reference-stop} gives the reference terminal scales.
On the top torus, \cref{thm:rgfl-terminal-zero-pair} selects the unique
controlled pair with both terminal couplings zero.
\Cref{thm:rgfl-common-pair} transfers this pair and its scalar bounds
to each eligible smaller torus. Its proof gives the stated local
agreement of both activity families, while
\cref{prop:rgfl-continuation} supplies their full norm bounds.
The initial mass estimate and its zero extension are
\eqref{eq:rgfl-bare-mass-bound}. The same common-pair proof establishes
independence from the smaller side, twist and external insertion.
\end{proof}
